\documentclass{amsart}
% For dealing with references we use the comment environment
\usepackage{verbatim}
\newenvironment{reference}{\comment}{\endcomment}
%\newenvironment{reference}{}{}
\newenvironment{slogan}{\comment}{\endcomment}
\newenvironment{history}{\comment}{\endcomment}

% For commutative diagrams we use Xy-pic
\usepackage[all]{xy}

% We use 2cell for 2-commutative diagrams.
\xyoption{2cell}
\UseAllTwocells

% We use multicol for the list of chapters between chapters
\usepackage{multicol}

% This is generally recommended for better output
\usepackage{lmodern}
\usepackage[T1]{fontenc}

% For cross-file-references

% Package for hypertext links:
\usepackage{hyperref}

% For any local file, say "hello.tex" you want to link to please

% Theorem environments.
%
\theoremstyle{plain}
\newtheorem{theorem}[subsection]{Theorem}
\newtheorem{proposition}[subsection]{Proposition}
\newtheorem{lemma}[subsection]{Lemma}

\theoremstyle{definition}
\newtheorem{definition}[subsection]{Definition}
\newtheorem{example}[subsection]{Example}
\newtheorem{exercise}[subsection]{Exercise}
\newtheorem{situation}[subsection]{Situation}

\theoremstyle{remark}
\newtheorem{remark}[subsection]{Remark}
\newtheorem{remarks}[subsection]{Remarks}

\numberwithin{equation}{subsection}

% Macros
%
\def\lim{\mathop{\mathrm{lim}}\nolimits}
\def\colim{\mathop{\mathrm{colim}}\nolimits}
\def\Spec{\mathop{\mathrm{Spec}}}
\def\Hom{\mathop{\mathrm{Hom}}\nolimits}
\def\Ext{\mathop{\mathrm{Ext}}\nolimits}
\def\SheafHom{\mathop{\mathcal{H}\!\mathit{om}}\nolimits}
\def\SheafExt{\mathop{\mathcal{E}\!\mathit{xt}}\nolimits}
\def\Sch{\mathit{Sch}}
\def\Mor{\mathop{\mathrm{Mor}}\nolimits}
\def\Ob{\mathop{\mathrm{Ob}}\nolimits}
\def\Sh{\mathop{\mathit{Sh}}\nolimits}
\def\NL{\mathop{N\!L}\nolimits}
\def\CH{\mathop{\mathrm{CH}}\nolimits}
\def\proetale{{pro\text{-}\acute{e}tale}}
\def\etale{{\acute{e}tale}}
\def\QCoh{\mathit{QCoh}}
\def\Ker{\mathop{\mathrm{Ker}}}
\def\Im{\mathop{\mathrm{Im}}}
\def\Coker{\mathop{\mathrm{Coker}}}
\def\Coim{\mathop{\mathrm{Coim}}}

% Boxtimes
%
\DeclareMathSymbol{\boxtimes}{\mathbin}{AMSa}{"02}

%
% Macros for moduli stacks/spaces
%
\def\QCohstack{\mathcal{QC}\!\mathit{oh}}
\def\Cohstack{\mathcal{C}\!\mathit{oh}}
\def\Spacesstack{\mathcal{S}\!\mathit{paces}}
\def\Quotfunctor{\mathrm{Quot}}
\def\Hilbfunctor{\mathrm{Hilb}}
\def\Curvesstack{\mathcal{C}\!\mathit{urves}}
\def\Polarizedstack{\mathcal{P}\!\mathit{olarized}}
\def\Complexesstack{\mathcal{C}\!\mathit{omplexes}}
% \Pic is the operator that assigns to X its picard group, usage \Pic(X)
% \Picardstack_{X/B} denotes the Picard stack of X over B
% \Picardfunctor_{X/B} denotes the Picard functor of X over B
\def\Pic{\mathop{\mathrm{Pic}}\nolimits}
\def\Picardstack{\mathcal{P}\!\mathit{ic}}
\def\Picardfunctor{\mathrm{Pic}}
\def\Deformationcategory{\mathcal{D}\!\mathit{ef}}

\setlength{\emergencystretch}{3em}
\begin{document}
\title[Stacks Project, Tag 074Q]{062 --- Skolem-Noether conjugacy theorem}
\maketitle
\noindent Copyright (C) 2005--2025 Johan de Jong. Stacks Project authors. Source: \href{https://stacks.math.columbia.edu/tag/074Q}{Tag 074Q}.

\noindent Editorial difficulty note (not author text). Difficulty H2: The author uses the structure of simple modules and their centralizers to compare two embeddings into a central simple algebra. An isomorphism of the resulting module structures is converted to inner conjugacy, with the module-structure lemmas retained..

\noindent Editorial provenance note (not author text). History: excerpt from revision \texttt{a0444\allowbreak 6e57e\allowbreak c1fbc\allowbreak 252a8\allowbreak 71afc\allowbreak ec775\allowbreak 2fb28\allowbreak 07b14}, brauer.tex, lines 483--508. Reference presentation was adapted; original-author context and core support blocks were retained or added. The mathematical wording is unchanged. Licensed under GNU FDL 1.2 or later; no invariant sections or cover texts. See ../licenses/COPYING. Context blocks reproduce author definitions and supporting results; they are not additional counted problems.

\begin{lemma}
\label{lemma-simple-module-unique}
Let $A$ be a finite simple $k$-algebra.
\begin{enumerate}
\item There exists exactly one simple $A$-module $M$ up to isomorphism.
\item Any finite $A$-module is a direct sum of copies of a simple module.
\item Two finite $A$-modules are isomorphic if and only if they
have the same dimension over $k$.
\item If $A = \text{Mat}(n \times n, K)$ with $K$ a finite skew field
extension of $k$, then $M = K^{\oplus n}$ is a simple $A$-module and
$\text{End}_A(M) = K^{op}$.
\item If $M$ is a simple $A$-module, then $L = \text{End}_A(M)$
is a skew field finite over $k$ acting on the left on $M$, we have
$A = \text{End}_L(M)$, and the centers of $A$ and $L$ agree.
Also $[A : k] [L : k] = \dim_k(M)^2$.
\item For a finite $A$-module $N$ the algebra $B = \text{End}_A(N)$ is a
matrix algebra over the skew field $L$ of (5). Moreover $\text{End}_B(N) = A$.
\end{enumerate}
\end{lemma}

\begin{proof}
By
Theorem \href{https://stacks.math.columbia.edu/tag/0747}{theorem wedderburn (Tag 0747)}
we can write $A = \text{Mat}(n \times n, K)$ for some finite skew
field extension $K$ of $k$. By
Lemma \href{https://stacks.math.columbia.edu/tag/074D}{lemma matrix algebras (Tag 074D)}
the category of modules over $A$ is equivalent to the category of
modules over $K$. Thus (1), (2), and (3) hold
because every module over $K$ is free. Part (4) holds
because the equivalence transforms the $K$-module $K$
to $M = K^{\oplus n}$. Using $M = K^{\oplus n}$ in (5)
we see that $L = K^{op}$. The statement about the center of $L = K^{op}$
follows from
Lemma \href{https://stacks.math.columbia.edu/tag/074D}{lemma matrix algebras (Tag 074D)}.
The statement about $\text{End}_L(M)$ follows from the explicit form
of $M$. The formula of dimensions is clear.
Part (6) follows as $N$ is isomorphic to a direct sum of
copies of a simple module.
\end{proof}

\begin{lemma}
\label{lemma-tensor-simple}
Let $A$, $A'$ be two simple $k$-algebras one of which is finite and central
over $k$. Then $A \otimes_k A'$ is simple.
\end{lemma}

\begin{proof}
Suppose that $A'$ is finite and central over $k$.
Write $A' = \text{Mat}(n \times n, K')$, see
Theorem \href{https://stacks.math.columbia.edu/tag/0747}{theorem wedderburn (Tag 0747)}.
Then the center of $K'$ is $k$ and we conclude that
$A \otimes_k K'$ is simple by
Lemma \ref{lemma-generate-two-sided-ideal}.
Hence $A \otimes_k A' = \text{Mat}(n \times n, A \otimes_k K')$ is simple
by Lemma \href{https://stacks.math.columbia.edu/tag/074D}{lemma matrix algebras (Tag 074D)}.
\end{proof}

\begin{lemma}
\label{lemma-generate-two-sided-ideal}
Let $A$ be a $k$-algebra. Let $K$ be a central $k$-algebra
which is a skew field. Then any two-sided ideal $I \subset A \otimes_k K$
is of the form $J \otimes_k K$ for some two-sided ideal $J \subset A$.
In particular, if $A$ is simple, then so is $A \otimes_k K$.
\end{lemma}

\begin{proof}
Set $J = \{a \in A \mid a \otimes 1 \in I\}$. This is a two-sided ideal
of $A$. And $I = J \otimes_k K$ by
Lemma \ref{lemma-generate-two-sided-sub}.
\end{proof}

\begin{lemma}
\label{lemma-generate-two-sided-sub}
Let $V$ be a $k$ vector space. Let $K$ be a central $k$-algebra
which is a skew field. Let $W \subset V \otimes_k K$ be a two-sided
$K$-sub vector space. Then $W$ is generated as a left $K$-vector
space by $W \cap (V \otimes 1)$.
\end{lemma}

\begin{proof}
Let $V' \subset V$ be the $k$-sub vector space generated by $v \in V$
such that $v \otimes 1 \in W$. Then $V' \otimes_k K \subset W$ and
we have
$$
W/(V' \otimes_k K)  \subset  (V/V') \otimes_k K.
$$
If $\overline{v} \in V/V'$ is a nonzero vector such that
$\overline{v} \otimes 1$ is contained in $W/(V' \otimes_k K)$,
then we see that $v \otimes 1 \in W$ where $v \in V$ lifts $\overline{v}$.
This contradicts our construction of $V'$. Hence we may replace
$V$ by $V/V'$ and $W$ by $W/(V' \otimes_k K)$ and it suffices to prove
that $W \cap (V \otimes 1)$ is nonzero if $W$ is nonzero.

\medskip\noindent
To see this let $w \in W$ be a nonzero element which can be written
as $w = \sum_{i = 1, \ldots, n} v_i \otimes k_i$ with $n$ minimal.
We may right multiply with $k_1^{-1}$ and assume that $k_1 = 1$.
If $n = 1$, then we win because $v_1 \otimes 1 \in W$.
If $n > 1$, then we see that for any $c \in K$
$$
c w - w c = \sum\nolimits_{i = 2, \ldots, n} v_i \otimes (c k_i - k_i c) \in W
$$
and hence $c k_i - k_i c = 0$ by minimality of $n$.
This implies that $k_i$ is in the center of $K$ which is $k$ by
assumption. Hence $w = (v_1 + \sum k_i v_i) \otimes 1$ contradicting
the minimality of $n$.
\end{proof}

\begin{theorem}
\label{theorem-skolem-noether}
Let $A$ be a finite central simple $k$-algebra. Let $B$ be a simple
$k$-algebra. Let $f, g : B \to A$ be two $k$-algebra homomorphisms.
Then there exists an invertible element $x \in A$ such that
$f(b) = xg(b)x^{-1}$ for all $b \in B$.
\end{theorem}

\begin{proof}
Choose a simple $A$-module $M$. Set $L = \text{End}_A(M)$.
Then $L$ is a skew field with center $k$ which acts on the left on $M$, see
Lemmas \href{https://stacks.math.columbia.edu/tag/0746}{lemma simple module (Tag 0746)} and \ref{lemma-simple-module-unique}.
Then $M$ has two $B \otimes_k L^{op}$-module structures defined by
$m \cdot_1 (b \otimes l) = lmf(b)$ and $m \cdot_2 (b \otimes l) = lmg(b)$.
The $k$-algebra $B \otimes_k L^{op}$ is simple by
Lemma \ref{lemma-tensor-simple}. Since $B$ is simple, the existence of a
$k$-algebra homomorphism $B \to A$ implies that $B$ is finite. Thus
$B \otimes_k L^{op}$ is finite simple and we conclude the two
$B \otimes_k L^{op}$-module structures on $M$
are isomorphic by Lemma \ref{lemma-simple-module-unique}.
Hence we find $\varphi : M \to M$ intertwining these operations.
In particular $\varphi$ is in the commutant of $L$ which implies that
$\varphi$ is multiplication by some $x \in A$, see
Lemma \ref{lemma-simple-module-unique}. Working out the definitions we see
that $x$ is a solution to our problem.
\end{proof}
\end{document}
